\documentstyle[%


righttag%


,11pt%


,amssymb%


,subeqn%


,russian%               remove in Eng.


,izvru%                 Set izven% in Eng.


]{amsart}





% 


%\newcommand{\r}[1]{\mbox{\rm{#1}}}


\newcommand{\vt}{\vartheta}


\newcommand{\La}{\Lambda}


\newcommand{\la}{\lambda}


\newcommand{\al}{\alpha}


\newcommand{\be}{\beta}


\newcommand{\om}{\omega}


\newcommand{\Om}{\Omega}


\newcommand{\ovo}{\ov{\omega}}


\newcommand{\si}{\sigma}


\newcommand{\di}{\displaystyle}


\newcommand{\no}{\nonumber}


\newcommand{\pa}{\di\partial}


\newcommand{\ra}{\rightarrow}


\newcommand{\eps}{\varepsilon}


\newcommand{\de}{\delta}


\newcommand{\vp}{\varphi}





\newcommand{\og}{\ov{\G}}


\newcommand{\Oh}[1] {{\cal O} \left(#1\right)}


\newcommand{\oga}{\ov{G}}


\newcommand{\os}{\ov{S}}


\newcommand{\od}{\ov{D}}


\newcommand{\eq}{\equiv}


\newcommand{\g}{\gamma}


\newcommand{\G}{{\mit{\Gamma}}}


\newcommand{\cu}{\cal{U}}


\newcommand{\re}{\vspace{-1mm}}


\newcommand{\pe}{\vspace{-2mm}}


\numberwithin{equation}{section}


\def\alph{\asbuk}  %% comment out in Eng.





\newcommand{\tts}[1]{}





\theoremstyle{plain}


\theorembodyfont{\it}


\newtheorem{thms}{\hskip\parindent Теорема}[section]





\theoremstyle{definition}


\theorembodyfont{\rm}


\newtheorem{notes}{\hskip\parindent Замечание}[section]





%\hoffset=-15mm%                  For Eng.


\setcounter{page}{73}


\god{2005}


\nomer{1~(512)} %                        Remove in Eng.


%\nomer{Vol.\,49, No.\,1}%               Set in Eng.


%\str{\pageref{begin}--\pageref{end}}%  Set in Eng.


\udk{519.632}





\author{\it Г.И.\,Шишкин}


% NB:   ^^ remove in Eng.


\title{О методе адаптирующихся сеток для сингулярно


возмущенных эллиптических уравнений реакции--диффузии в области


с криволинейной границей}


\thanks{Работа выполнена при


финансовой поддержке Российского фонда фундаментальных


исследований (коды проектов \No\No\,04-01-00578, 04-01-89007-НВО\_a)


и Нидерландской


организации научных исследований NWO (проект \No\,047.016.008).}





\date{}


%\pagestyle{myheadings}%         Set in Eng.


\pagestyle{plain} %              Remove in Eng.


\begin{document}


%\thispagestyle{plain}%          Set in Eng.


\maketitle


\label{begin}








В области с криволинейной границей рассматривается краевая задача


для сингулярно возмущенных эллиптических уравнений


реакции--диффузии. Монотонные классические разностные схемы для


задач такого типа сходятся лишь при условии $\eps\gg


N_1^{-1}+N_2^{-1}$, где $\eps$ --- возмущающий параметр, величины


$N_1$ и $N_2$ определяют число узлов сетки по переменным $x_1$ и


$x_2$. Поэтому для сингулярно возмущенных уравнений в области с


криволинейной границей необходимо развивать специальные численные


методы, ошибки решений которых довольно слабо зависят от параметра


$\eps$ и, в частности, не зависят от $\eps$  (т.\,е.


$\eps$-равномерно сходящиеся методы). Исследуются


схемы на адаптивных сетках, локально сгущающихся в окрестности


погранслоя. Оказывается, что в классе разностных схем на основе


классических аппроксимаций задачи на прямоугольных сетках,


априорно, либо апостериорно локально сгущающихся в погранслое, не


существует схем, сходящихся $\eps$-равномерно и даже при условии


$\eps\approx N_1^{-2}+N_2^{-2}$, если общее число узлов локально


сгущающейся сетки порядка $N_1N_2$. Таким образом,


непосредственное использование технологии на основе адаптивных


сеток не позволяет существенно расширить область сходимости


классических численных методов. Использование техники поперечников по


Колмогорову позволило установить условия, необходимые для


$\eps$-равномерной сходимости (при $P\ra\infty$, где $P$ --- общее число


узлов сетки) аппроксимаций решений краевых задач; требования,


вытекающие из этих условий, кладутся в основу построения


специальных сеток. Использование таких сгущающихся (в слое) сеток,


однако, в локальной системе координат, адаптирующейся к границе


области, позволяет строить схемы, сходящиеся $\eps$-равномерно


при $P\ra\infty$.








\section{Введение}


\label{section.1}





Нередко при исследовании процессов тепломассообмена в средах с


малыми коэффициентами теплопроводности/диффузии возникают краевые


задачи для сингулярно возмущенных уравнений в частных производных


(с возмущающим параметром $\eps$ --- коэффициентом при старших


производных уравнений), в частности, для эллиптических уравнений


реакции--диффузии в областях с криволинейными границами. Решения


таких задач при малых значениях параметра $\eps$ имеют особенности


типа пограничных слоев. Эти особенности в сочетании с


криволинейностью границ вызывают трудности при численном решении


задачи (классическая разностная схема сходится лишь при условии


$\eps\gg N_1^{-1}+N_2^{-1}$, см., напр., утверждение теоремы


\ref{t3.1}). Для такого типа задачи в \cite{1} построена


специальная разностная (итерационная) схема, сходящаяся


$\eps$--равномерно со скоростью $\Oh{N^{-1}\ln N+q^{-n}}$, где


$N=\min[N_1,N_2]$, $q<1$, $n$ --- номер итерации. При построении


схемы в окрестности границы области использовались сгущающиеся в


погранслое сетки, являющиеся прямоугольными в координатах,


связанных с границей. При построении таких сеток, согласованных с


границей, используются криволинейные координаты, в которых граница


области есть координатная линия.


 В связи с достаточной сложностью этой схемы


возникает интерес к альтернативным методам построения специальных


схем.





В случае сингулярно возмущенных задач в областях с криволинейными


границами при разработке специальных численных методов,


погрешность решений которых достаточно слабо зависит от величины


параметра $\eps$ и, в частности, $\eps$-равномерно сходящихся


методов, представляется целесообразным использовать методы на


основе локально сгущающихся сеток. Такие методы достаточно хорошо


разработаны для ряда сингулярно возмущенных краевых задач на


прямоугольных областях (см., напр., \cite{1}--\cite{4} и


библиографию там же). Так, в случае регулярных краевых задач,


решения которых имеют особенности (большие производные), эффект


повышения точности приближенного решения может быть достигнут за


счет априорного и/или апостериорного локального сгущения


прямоугольной сетки в тех подобластях, где ошибки приближенного


решения большие (напр., \cite{5}--\cite{7}).


В случае сингулярно возмущенных уравнений в областях с криволинейными


границами представляют интерес численные методы на основе (достаточно


простых) сеток, локально сгущающихся в пограничном слое, однако не


использующих криволинейные координаты, связанные с границей области (или,


короче, на основе локально сгущающихся сеток, не согласованных в погранслое


с границей).





В данной работе показано, что непосредственное использование


такого подхода в случае сингулярно возмущенных задач в областях с


криволинейными границами не~является достаточно эффективным. Для


классических аппроксимаций задачи реакции--диффузии на адаптирующихся,


локально сгущающихся сетках (прямоугольных в окрестности погранслоя)


сходимость сеточных решений недостижима уже при условии


$\eps^{1/2}\approx N_1^{-1}+N_2^{-1}$ (см., напр.,


утверждение теоремы \ref{t4.1}). Использование техники поперечников по


Колмогорову позволяет для распределения сеточных узлов указать


условия, необходимые, а также и достаточные, для


$\eps$-равномерной сходимости (с ростом числа узлов сетки,


определяющей триангуляцию области) аппроксимаций


решений краевой задачи (см. раздел~\ref{section.5}). В частности,


при построении схем, сходящихся, например, при $N_1^{-1}+N_2^{-1}=


\Oh{\eps^{\alpha}}$, $\alpha<4^{-1}$, требуется использовать


(сгущающиеся в погранслое) сетки на основе криволинейной системы


координат, согласованной с границей области. В то же время локально


сгущающиеся сетки, согласованные в погранслое с границей,


позволяют построить схему, сходящуюся


$\eps$-равномерно со скоростью


$\cal O(N_1^{-2}\ln^2N_1+N_2^{-2}\ln^2N_2)$ (пример построения


специальной схемы см. в разделе~\ref{section.6}).





Заметим, что в \cite{12} при построении специальных схем для


рассматриваемого класса задач применялись локально сгущающиеся


сетки, порождаемые прямыми, параллельными координатным осям.


Схемы на таких сетках сходятся на $\ov{D}$ --- множестве


определения решения --- при более слабом условии, накладываемом на


параметр $\eps$, чем в случае классических разностных схем,


а также сходятся $\eps$-равномерно, однако лишь на подмножествах из


$\ov{D}$, а именно, вне пограничного слоя и


на достаточно узких множествах, ``пересекающих'' пограничный слой.


Как следует из данной работы, достаточно простые специальные сетки,


используемые в \cite{12}, не позволяют строить схемы, сходящиеся


$\eps$-равномерно на всем множестве $\ov{D}$.








\section{Постановка задачи. Цель исследования}


\label{section.2}








{\ref{section.2}.1.} В ограниченной области $D$ с достаточно


гладкой границей $\varGamma=\od\setminus D$ рассмотрим задачу Дирихле


для сингулярно возмущенного эллиптического уравнения


реакции--диффузии


\begin{subequations} \label{eq.2.1}


\begin{alignat}{2} \label{eq.2.1a}


Lu(x) & = f(x), &\quad x&\in D; \\


\label{eq.2.1b}


u(x) & = \vp(x), &\quad x&\in\varGamma.


\end{alignat}


\end{subequations}


Здесь


$$


L \equiv \eps^2 L_2+L_0,\quad L_0\equiv-c_0(x),\quad


L_2\equiv \sum_{s,k=1}^2 a_{sk}(x)


\frac{\pa^2}{\pa x_{ s}\pa x_{ k}}+\sum_{s=1}^2b_s(x)


\frac{\pa}{\pa x_s}-c(x);


$$


коэффициенты и правая часть уравнения, а также граничная функция


предполагаются достаточно гладкими и удовлетворяющими


условию\footnote{Здесь и ниже через $M$ (через $m$) обозначаем


достаточно большие (малые) положительные постоянные, не зависящие


от величины параметра $\eps$ и от параметров разностных схем.


Запись $L_{(j.k)}$ ($M_{(j.k)}$, $D_{h(j.k)}$) означает, что эти


операторы (постоянные, сетки) введены в формуле $(j.k)$.}


\begin{gather} \label{eq.2.2}


a_0(\xi_1^2+\xi_2^2)\leq \sum_{s,k=1}^2


a_{sk}(x)\xi_s \xi_k \leq a^0(\xi_1^2+\xi_2^2), \ \


c_0(x)\geq c_{00},\ \ c(x)\geq 0, \\


|f(x)|\leq M, \ \ x\in \od; \quad


|\vp(x)|\leq M, \ \ x\in\varGamma, \quad a_0,c_{00}>0; \notag


\end{gather}


параметр $\eps$ принимает произвольные значения из полуинтервала $(0,1]$.





При $\eps\ra 0$ в окрестности границы $\varGamma$ появляется (регулярный)


пограничный слой, экспоненциально убывающий при удалении от $\varGamma$.


\medskip





{\ref{section.2}.2.} Ошибки решений разностных схем, построенных


на основе классических разностных аппроксимаций задачи


(\ref{eq.2.1}), зависят от величины параметра $\eps$ и становятся


малыми лишь при значениях параметра $\eps$, существенно


превосходящих ``эффективные'' шаги сеток по $x_1$, $x_2$. В силу


оценки (\ref{eq.3.10}) классическая разностная схема


(\ref{eq.3.4}), (\ref{eq.3.9}) (см. раздел~\ref{section.3})


сходится при условии


\begin{equation}\label{eq.2.3}


\eps\gg N_1^{-1}+N_2^{-1},


\end{equation}


где величины $N_1$, $N_2$ определяют число узлов сеток по $x_1$,


$x_2$. При нарушении этого условия решения разностной схемы не


сходятся к решению задачи (\ref{eq.2.1}).





В связи с отмеченным поведением сеточных решений возникает интерес


к построению специальных разностных схем, погрешность решений


которых не зависит от величины параметра $\eps$. В частности,


представляют интерес разностные схемы, сходящиеся при более слабом


условии, чем (\ref{eq.2.3}) --- условие сходимости решений


сеточных задач (\ref{eq.3.4}), (\ref{eq.3.7}) и


(\ref{eq.3.4}), (\ref{eq.3.9}).





Цель работы --- для краевой задачи (\ref{eq.2.1}) с использованием


метода сгущающихся сеток построить разностные схемы, сходящиеся


$\eps$-равномерно, а также схемы, близкие к таковым, а именно,


сходящиеся при значениях $\eps$, много меньших, чем в условии


(\ref{eq.2.3}). Указываются условия, накладываемые на


распределение сеточных узлов, которые обеспечивают построение схем


указанного типа.








\section{Априорные оценки}


\label{section.7}





Приведем априорные оценки решения задачи (\ref{eq.2.1}),


используемые в дальнейших построениях (см. также \cite{1}--\cite{4}).





С использованием принципа максимума, а также априорных оценок для


регулярных краевых задач \cite{13} (внутренних оценок и оценок вплоть


до гладких частей границы) устанавливаются оценки


\begin{subequations} \label{eq.7.1.1}


\begin{alignat}{2}\label{eq.7.1.1a}


|u(x)&|\le M, &\quad \ x&\in\od;\\


\label{eq.7.1.1b}


\bigg|\frac{\pa^k}{\pa x_1^{k_1}\pa x_2^{k_2}}\,u(x)\bigg|&\le M


\eps^{-k}, &\quad x&\in \od,\ \ k\le K;


\end{alignat}


\end{subequations}


при достаточно гладких данных задачи величина $K$ может быть выбрана достаточно


большой. Заметим, что в переменных $\zeta=(\zeta_1,\zeta_2)$,


$\zeta_s=\eps^{-1}x_s$, $s=1,2$, дифференциальное уравнение является


регулярным.





При выводе оценок на основе асимптотики решение


задачи представим в виде суммы функций


\begin{equation}\label{eq.7.1}


u(x)=U(x)+V(x),\quad x\in\od,


\end{equation}


где $U(x)$ и $V(x)$ --- регулярная и сингулярная части решения.


Функция $U(x)$ --- сужение на $\od$ функции $U^0(x)$, $x\in R^2$,


являющейся ограниченным решением уравнения


$$


L^0U^0(x)=f^0(x), \quad x\in R^2;


$$


оператор $L^0$ и функция $f^0(x)$ --- продолжения оператора


$L_{(\ref{eq.2.1})}$ и функции $f_{(\ref{eq.2.1})}(x)$ с множества


$\od$ на $R^2$ с сохранением свойств. Для простоты вне ближайшей окрестности


множества $\od$ функцию $f^0(x)$ считаем равной нулю, а оператор


$L^0$ определяем соотношением $L^0\equiv \eps^2\Delta -1$.


Функция $V(x)$, $x\in\od$, --- решение задачи


\begin{alignat*}{2}


LV(x) &= 0, &\quad x&\in D,\\


V(x) &= \vp(x)-U(x), &\quad x&\in\varGamma.


\end{alignat*}





Функцию $U^0(x)$, $x\in R^{2}$, представим в виде суммы функций


$$


U^0(x)=\sum_{k=0}^{n}\,\eps^{2k}\,U^0_k(x)+v_{U}^n(x) \eq


U^{0n}(x)+v_{U}^n(x), \quad x \in R^{2},


$$


где $n \ge 0$; $v_{U}^n(x)$ -- остаточный член разложения. Функции


$U^0_k(x)$ --- компоненты регулярной части решения являются решениями задач


\begin{gather*}


L_0^0\,U^0_0(x) \eq - c_0^{0}(x) \,U^0_0(x)=f^0(x), \quad 


x \in R^{2}; \\


L_0^0U^0_k(x) = -L_2^0\,U_{k-1}^0(x), \quad 


x \in R^{2}, \quad 0<k\le n,


\end{gather*}


где


$$


L_2^0\equiv \sum_{s,k =1}^2 a_{sk}^0(x)\frac{\partial^2}{\partial


x_s\partial x_k}+ \sum_{s=1}^2 b_s^0(x)\frac{\partial}{\partial x_s}-


c^0(x),


$$


функции $a^{0}_{sk}(x),\dotsc,c^{0}_0(x)$ --- продолжения


функций $a_{sk}(x),\dotsc,c_0(x)$, определяющих оператор


$L$. С учетом оценок компонент $U_k^0(x)$, $\,v^n_U(x)$ находится оценка для


производных функции $U^0(x)$, $x\in\od$ (остаточный член $v^n_U(x)$


оцениваем с учетом оценки (\ref{eq.7.1.1b}), см., напр., \cite{2}).


Для функции $U(x)$ получаем оценку


$$


\bigg|\frac{\pa^k}{\pa x_1^{k_1}\pa x_2^{k_2}}\,U(x)\bigg|\le M


[1+\eps^{2(n+1)-k}], \quad x\in \od, \ \ k\le K.


$$


При $n\ge 2^{-1}K-1$ имеем


\begin{subequations}\label{eq.7.2}


\begin{equation}\label{eq.7.2a}


\bigg|\frac{\pa^k}{\pa x_{ 1}^{k_1}\pa x_{ 2}^{k_2}}\,U(x)


\bigg|\leq M,\quad x\in\od,\quad k\leq K.


\end{equation}


Для функции $V(x)$ с учетом оценки (\ref{eq.7.2a}) находим


\begin{equation}\label{eq.7.2b}


\bigg|\frac{\pa^k}{\pa x_{ 1}^{k_1}\pa x_{ 2}^{k_2}}V(x)


\bigg|\leq M\eps^{-k}\exp(-m\eps^{-1}r(x,\varGamma)),


\quad \ \ x\in\od,\quad k\leq K,


\end{equation}


\end{subequations}


где $r(x,\varGamma)$ --- расстояние от точки $x$ до множества $\varGamma$.





Уточним оценку (\ref{eq.7.2b}). Область $\od$ представим в виде


суммы перекрывающихся подобластей


\begin{equation}\label{eq.7.3}


\od=\od{}^1\cup\od{}^2,


\end{equation}


где множество $\od{}^1$ есть ``кольцо'' --- $m_1$-окрестность границы


$\varGamma$, а множество $\od{}^2$ не имеет общих точек с $\varGamma$. Считаем,


что $r(x,\varGamma)\leq m_1$, $x\in\od{}^1$, где постоянная $m_1$ меньше


радиуса кривизны границы $\varGamma$; $r(x^2,\varGamma)\geq m$, $x^2\in\od{}^2$;


$r(x^1,x^2)\geq\de$, $x^i\in\od{}^i$, $i=1,2$,


$x^1,x^2\notin D^1\cap D^2$, $\,\de$ --- ширина


минимального перекрытия подобластей.





На множестве $\od{}^1$ перейдем к переменным $\xi=(\xi_1,\xi_2)$:


$$


\xi_s=\xi_s(x),\quad s=1,2,


$$


где $\xi_1(x)$ --- расстояние от точки $x$ до границы $\varGamma$,


$\xi_2(x)$ --- длина дуги границы $\varGamma$ (при обходе ее в заданном


направлении) от фиксированной точки $x^0\in\varGamma$ до точки


$x^1\in\varGamma$, ближайшей к точке $x$. Для функций $v(x)$,


$x\in\od{}^1$, $W(\xi)$, $\xi\in\ov{\wt{D}}{}^1$, где


$\wt{D}{}^1=\{\xi:\xi=\xi(x)$, $x\in D^1\}$, будем использовать


обозначения


\begin{gather*}


v(x(\xi))=v_{\xi}(\xi)=\{v(x)\}_{\xi},\quad W(\xi(x))=W_x(x), \\


D^0_{\xi}=\xi(D^0)=\{\xi: \xi=\xi(x),\ x\in D^0\},


\quad D^0\subseteq\od{}^1. 


\end{gather*}


Здесь $x=x(\xi)$ --- отображение, обратное $\xi=\xi(x)$. Полагаем


$$


\wt{D}{}^0_x=x(\wt{D}{}^0)=\{x: x=x(\xi),\ \xi\in\wt{D}^0 \},


$$


где $\wt{D}{}^0$ --- некоторое подмножество из $\ov{\wt{D}}{}^1$.





В новых переменных уравнение (\ref{eq.2.1a}) на $\wt{D}{}^1$ принимает вид


$$


\wt{L}U(\xi)=F(\xi),\quad \xi\in\wt{D}{}^1.


$$


Здесь


$$


\wt{L}=L_{\xi}\equiv \eps^2\bigg\{\sum_{s,k=1}^2 A_{sk}(\xi)


\frac{\pa^2}{\pa \xi_{ s}\pa \xi_{ k}}+\sum_{s=1}^2 B_s(\xi)


\frac{\pa}{\pa \xi_s}-C(\xi) \bigg\}-C_0(\xi),


$$


коэффициенты оператора $\wt{L}$ определяются через коэффициенты


оператора $L_{(\ref{eq.2.1})}$:


\begin{alignat*}{2}


A_{sk}(\xi)&=\bigg\{


\sum_{r,p=1,2} a_{rp}(x)\frac{\partial}{\partial x_r} \xi_{s}(x)


\frac{\partial}{\partial x_p} \xi_{k}(x)\bigg\}_{\xi}, &\quad s,k&=1,2;\\


B_{s}(\xi)&=\bigg\{


\sum_{r=1,2} b_{r}(x)\frac{\partial}{\partial x_r} \xi_{s}(x)+


\sum_{r,p=1,2} a_{rp}(x)\frac{\partial^2}{\partial x_r\partial x_p}


\xi_s(x)\bigg\}_{\xi}, &\quad s&=1,2,


\end{alignat*}


функции $C$, $C_0$, $F$, $U$ определяются


соотношением $V(\xi)=v_{\xi}(\xi)$.





Задаче (\ref{eq.2.1}) в случае разбиения области (\ref{eq.7.3})


соответствует задача


\begin{subequations} \label{eq.7.4}


\begin{align}\label{eq.7.4a}


\wt{L} U(\xi) &= F(\xi), \quad


\xi \in \wt{D}^1, \\


U(\xi) &=\begin{cases}


\Phi(\xi), & \xi\in\wt{\varGamma}^1\cap \wt{\varGamma}; \\


u_{\xi}(\xi), & \xi\in\wt{\varGamma}^1\setminus \wt{\varGamma},


\end{cases}


\notag\\


%


\label{eq.7.4b}


L u(x) &=f(x), \quad x\in D^2, \\


u(x)& = U_x(x), \quad x\in\varGamma^2,\notag


\end{align}


\end{subequations}


где $\Phi(\xi)=\vp_{\xi}(\xi)$. Подзадача (\ref{eq.7.4a}) есть


периодическая задача (по переменной $\xi_2$) на (вертикальной)


полосе.





Оценим производную $(\pa/\pa\xi_2)U^S(\xi)\equiv U_1^S(\xi)$ сингулярной


компоненты $U^S(\xi)$ в представлении


\begin{equation}\label{eq.7.5}


U(\xi)=U^R(\xi)+U^S(\xi),\quad \xi\in \ov{\wt{D}}{}^1;


\end{equation}


здесь $U^R(\xi)$ и $U^S(\xi)$ --- регулярная и сингулярная части


решения задачи (\ref{eq.7.4a}),


$U^R(\xi)=\{U(x)\}_{\xi}$, $U^S(\xi)=\{V(x)\}_{\xi}$,


$U(x)$ и $V(x)$ --- компоненты из (\ref{eq.7.1}). В силу оценки


(\ref{eq.7.2a}) функция $U_1^S(\xi)$ ограничена $\eps$-равномерно на


множестве $\wt{\varGamma}^1\cap \wt{\varGamma}$. Дифференцируя по


$\xi_2$ уравнение


$$


\wt{L} U^S(\xi)=0, \quad \xi\in \wt{D}{}^1,


$$


и принимая во внимание оценку (\ref{eq.7.2b}), находим


$$


\bigg|\frac{\pa}{\pa \xi_{2}}U^S(\xi)\bigg|\leq


M\exp(-m\eps^{-1}r(\xi,\wt{\varGamma})),\quad \xi\in \ov{\wt{D}}{}^1.


$$


Подобным образом оцениваем производные


$\big(\partial^k/\pa \xi_{ 1}^{k_1}\pa \xi_{ 2}^{k_2}\big)U^S(\xi)$,


$~\xi\in \ov{\wt{D}}{}^1$.





Для компонент из (\ref{eq.7.5}) получаются оценки


\begin{align}\label{eq.7.6}


\bigg|\frac{\pa^k}{\pa \xi_{ 1}^{k_1}\pa \xi_{ 2}^{k_2}}U^R(\xi)


\bigg|&\leq M, \\


\bigg|\frac{\pa^k}{\pa \xi_{ 1}^{k_1}\pa \xi_{ 2}^{k_2}}\,U^S(\xi)


\bigg|&\leq M\eps^{-k_1}\exp(-m\eps^{-1}


r(\xi,\wt{\varGamma})),\quad \xi\in \ov{\wt{D}}{}^1,\quad k\leq K. \notag


\end{align}


В оценках (\ref{eq.7.2}), (\ref{eq.7.6}) величина $m$ --- произвольная


постоянная из интервала $(0,m_0)$,


$m_0={\min\limits_{\od{}^1}}^{1/2}[(a^0)^{-1}c_0(x)]$,


$a^0=a^0_{(\ref{eq.2.3})}$.





Таким образом, установлена следующая





\begin{thms}\label{t7.1}


Пусть $a_{sk}, b_s, c, c_0, f\in C^{l+\al}(\od)$, $\vp\in


C^{l+\al}(\varGamma)$, $\varGamma\in C^{l+\al}$, $l\geq0$, $\al>0$, $s,k=1,2$.


Тогда для решений задачи $(\ref{eq.2.1})$ и его компонент из


представлений $(\ref{eq.7.1})$, $(\ref{eq.7.5})$ справедливы


оценки $(\ref{eq.7.1.1})$, $(\ref{eq.7.2})$, $(\ref{eq.7.6})$, где


$K=2[(l+2)/4];$ $[\,\cdot\,]$ --- целая часть числа.


\end{thms}








\section{Классические разностные схемы}


\label{section.3}





Выпишем классическую разностную схему для задачи (\ref{eq.2.1}) и


укажем некоторые трудности, возникающие при численном решении при


малых значениях параметра $\eps$.





На плоскости $R^2$ введем прямоугольные базовые сетки, на основе


которых будем строить сетки на $\od$. Пусть


\begin{equation}\label{eq.3.1}


R_h^2=\om_1\times\om_2,


\end{equation}


где $\om_s$ --- сетка на оси $x_s$, $s=1,2$; $\om_s$ --- сетки с


произвольным распределением узлов, удовлетворяющим лишь условию


$h\leq MN^{-1}$, где $h=\max\limits_s h_s$, $h_s=\max\limits_i h_s^i$,


$h_s^i=x_s^{i+1}-x_s^i$, $x_s^i,~x_s^{i+1}\in\om_s$,


$N=\min[N_1,N_2]$, $N_s+1$ --- минимальное число узлов сетки


$\om_s$ на отрезке единичной длины на оси $x_s$. Наибольший


интерес для нас будут представлять сетки, равномерные по


$x_1$ и $x_2$


\begin{equation}\label{eq.3.2}


R_h^2=R_{h(\ref{eq.3.1})}^2,


\end{equation}


где $\om_s$ --- равномерные сетки с шагами $h_s=N_s^{-1}$,


$s=1,2$. На множестве $\od$ строим сетку


$\od_h=\od_h(R_{h(\ref{eq.3.1})}^2)$ на основе базовой сетки


$R_{h(\ref{eq.3.1})}^2$ (напр., \cite{8}). Пусть


$D_h=D\cap R_h^2$ --- множество внутренних узлов и $\varGamma_h$


--- множество граничных узлов, есть множество точек пересечения


прямых $x_s=x_s^{i_s}$, $s=1,2$, с границей $\varGamma$,


$(x_1^{i_1},x_2^{i_2})\in D_h$. Полагаем


\begin{equation}\label{eq.3.3}


\od_h=D_h\cup\varGamma_h;


\end{equation}


для каждого узла $x\in D_h$ имеются ближайшие соседние узлы из


$\od_h$, образующие пятиточечный шаблон типа ``прямой крест'' с


центром в точке $x$. К строго внутренним узлам


$D_h^{(in)}=D_{h(\ref{eq.3.3})}^{(in)}$ отнесем те узлы $x$, для


которых все узлы девятиточечного шаблона типа ``ящик'' (с центром


$x$) принадлежат $D_h$. Заметим, что величина $P$ --- число узлов


множества $\od_h$, порядка $N_1N_2$.





Оператору $L_{(\ref{eq.2.1})}$ на сетке $\od_{h(\ref{eq.3.3})}$


сопоставим сеточный оператор


\begin{multline} \label{eq.3.3^1}


\La=\La(L;\od_h)\equiv\eps^2 \bigg \{\sum_{s=1}^{2}a_{ss}(x)


\de_{\ov{xs}\,\wh{x s}}+2^{-1}\eta(x)\sum\begin{Sb} s,k=1 \\


s\neq k\end{Sb}^{2}


[a_{sk}^{+}(x)(\de_{xs\,xk}+\de_{\ov{xs}\,\ov{xk}})+ \\


%


+a_{sk}^{-}(x)(\de_{xs\,\ov{xk}}+\de_{\ov{xs}\,xk})]+


\sum_{s=1}^{2}b_s(x)\de_{\wt{xs}}-c(x) \bigg \}-c_0(x). 


\end{multline}


Здесь $\de_{\ov{xs}\,\wh{x s}} z(x)=z_{\ov{xs}\,\wh{x s}}(x),


\dotsc,\de_{\ov{xs}\,xk} z(x)=z_{\ov{xs}\,xk}(x)$ и


$\de_{\wt{xs}}z(x)=z_{\wt{xs}}(x)$ --- вторые и первые


(центральные) разностные производные, например,


$\de_{\ov{x1}\,\wh{x1}}z(x)=2(h_1^i+h_1^{i-1})^{-1}


[\de_{x1}-\de_{\ov{x1}}]z(x)$, 


$\de_{\ov{x1}\,x2}\,z(x)= (h_1^{i-1} )^{-1}


[\de_{x2}\,z(x)-\de_{x2}\,z(x_1^{i-1},x_2)]$, 


$\de_{\ov{x1}}z(x)= (h_1^{i-1})^{-1}[z(x)-z(x_1^{i-1},x_2)]$,


$x=(x_1^i,x_2)$, $\eta(x)=1$,


$x\in D_h^{(in)}$, $\eta(x)=0$, $x\in D_h\setminus D_h^{(in)}$;


$v^{+}(x)=2^{-1}(v(x)+|v(x)|)$, $v^{-}(x)=2^{-1}(v(x)-|v(x)|)$.





Задачу (\ref{eq.2.1}) аппроксимируем разностной схемой


\begin{equation}


\begin{split}


\La z(x)&=f(x), \quad x\in D_h, \label{eq.3.4}\\


z(x)&=\vp(x), \quad x\in \varGamma_h.


\end{split}


\end{equation}





Разностная схема (\ref{eq.3.4}), (\ref{eq.3.3})


(при $N\geq N_0$, где $N_0$ --- достаточно большое число)


монотонна \cite{8}, если выполняется условие


\begin{equation}


\label{eq.3.5}


a_{12}\,(x)\equiv 0,\quad x\in\od.


\end{equation}


При нарушении этого условия непосредственно убеждаемся, что


схема (\ref{eq.3.4}), (\ref{eq.3.3})


монотонна, если для распределения узлов $\od_h$ выполняется условие


\begin{equation}


\label{eq.3.6}


\min_{s,D_h^{(in)}}\{2(h_s^{i_s}+h_s^{i_s-1})^{-1}a_{ss}(x)-


\max[(h_{3-s}^{i_{3-s}})^{-1},(h_{3-s}^{i_{3-s}-1})^{-1}]


|a_{12}(x)|-|b_s(x)|\}\geq 0;


\end{equation}


признаки монотонности см., например, в \cite{8}.


Это условие справедливо, например, если максимальный и минимальный


шаги сетки $\om_s$ соизмеримы, величины $N_1$ и $N_2$ одного


порядка, а коэффициент $a_{12}\,(x)$ мал по сравнению с


коэффициентами $a_{11}(x)$, $a_{22}(x)$.





Будем предполагать, что выполняется условие (\ref{eq.3.5})


либо~(\ref{eq.3.6}).





Принимая во внимание априорные оценки (\ref{eq.7.2}), (\ref{eq.7.6}),


с использованием техники мажорантных функций в случае сетки


\begin{equation}\label{eq.3.7}


\od_h=\od_h(R^2_{h(\ref{eq.3.1})})


\end{equation}


устанавливаем оценку


$$


%% \label{eq.3.8}


|u(x)-z(x)|\leq M(\eps+N^{-1})^{-1}N^{-1}, \quad x\in\od_h.


$$


Здесь в качестве мажорантной используем функцию $w(x)=\mathrm{const}$,


$x\in\ov{D}$. На сетке


\begin{equation}\label{eq.3.9}


\od_h=\od_h(R^2_{h(\ref{eq.3.2})})


\end{equation}


получается оценка


\begin{equation}\label{eq.3.10}


|u(x)-z(x)|\leq M\sum_{s=1,2}(\eps+N_s^{-1})^{-2}N_s^{-2},


\quad x\in\od_h.


\end{equation}








Введем некоторые определения.


Пусть на множестве $E_{N,\eps}=E_N\times E_{\eps}$, где


$E_N$ --- подмножество из множества пар натуральных чисел $N_1$, $N_2$,


удовлетворяющих условию


$N_1, N_2\ge M_0$, $E_\eps=\{\eps:\eps\in(0,1]\}$, определены функции


$\psi_i(N_1^{-1},N_2^{-1},\eps)$, $i=1,2$; 


$\psi_i(N_1^{-1},N_2^{-1},\eps)>0$. Запись


$\psi_1(N_1^{-1},N_2^{-1},\eps)=\wh{o}\big(\psi_2(N_1^{-1},N_2^{-1},\eps)


\big)$ на $E_{N,\eps}$ означает, что найдется такая точка


$(\wt{N}_1^{-1},\wt{N}_2^{-1},\wt{\eps})$,


для которой выполняется соотношение


\begin{gather*}


\psi_1(N_1^{-1},N_2^{-1},\eps)[\psi_2(N_1^{-1},N_2^{-1},\eps)]^{-1}


\to 0 \ \ \text{при} \ \ (N_1^{-1},N_2^{-1},\eps)\to


(\wt{N}_1^{-1},\wt{N}_2^{-1},\wt{\eps}),\\


(N_1^{-1},N_2^{-1},\eps),\, (\wt{N}_1^{-1},\wt{N}_2^{-1},\wt{\eps})


\in E_{N,\eps}.


\end{gather*}


Пусть для сеточной функции $z(x)$, $x\in\ov{D}_h$, ---


решения некоторой разностной схемы, выполняется оценка


$$


|u(x)-z(x)|\le M\mu(N_1^{-1},N_2^{-1},\eps), \quad 


x\in\ov{D}_h.


$$


Будем говорить, что эта {\it оценка неулучшаема} по вхождению


величин $N_1$, $N_2$, $\eps$, если оценка


$$


|u(x)-z(x)|\le M \mu_0 (N_1^{-1}, N_2^{-1},\eps),\ \ x \in \od_h,


$$


вообще говоря, неверна в том случае, когда $\mu_0 (N_1^{-1},


N_2^{-1},\eps) = \wh{o} (\mu(N_1^{-1}, N_2^{-1},\eps))$


на $E_{N,\eps}$.





Пусть при $N_1, N_2\to\infty$, $\eps\in E_{\eps}$, решение


разностной схемы сходится к решению краевой задачи в случае


условия $N_1^{-1},N_2^{-1}=o(\eps^{\nu})$, $\eps\in E_\eps$, однако 


сходимость при условии $N_1^{-1},N_2^{-1}=\Oh{\eps^{\nu}}$, вообще 


говоря, не имеет места. В этом


случае будем говорить, что разностная схема {\it сходится с


дефектом} $\nu$ относительно параметра $\eps$ (или, короче,


сходится с дефектом $\nu$). При $\nu=0$ сходимость схемы


$\eps$-равномерная.





Рассматривая решения модельных задач, убеждаемся, что оценка (\ref{eq.3.10})


является неулучшаемой по вхождению величин $N_1$, $N_2$, $\eps$.


Отсюда вытекает, что условие


$(\eps^{-1}N_1^{-1},\eps^{-1}N_2^{-1}\to 0$ при $N_1,N_2\to\infty$,


$(N_1^{-1},N_2^{-1},\eps)\in E_{N,\eps})$


\begin{equation}


\label{eq.3.11}


N_1^{-1},\, N_2^{-1}=o(\eps), \quad N_1,N_2\to\infty,


\end{equation}


является необходимым и достаточным для сходимости решений разностной схемы


(\ref{eq.3.4}), (\ref{eq.3.9}); схема сходится с дефектом $\nu=1$. %%%\vspace{1mm}





Таким образом, справедлива





\begin{thms} \label{t3.1}


Пусть для компонент решений краевой задачи $(\ref{eq.2.1})$ из


представлений $(\ref{eq.7.1})$, $(\ref{eq.7.5})$ выполняются


априорные оценки $(\ref{eq.7.2})$, $(\ref{eq.7.6})$, где $K=4$, и


пусть выполняeтся условие $(\ref{eq.3.5})$ либо $(\ref{eq.3.6})$.


Тогда решение разностной схемы $(\ref{eq.3.4})$, $(\ref{eq.3.9})$


сходится при условии $(\ref{eq.3.11});$


для сеточных решений справедлива оценка $(\ref{eq.3.10})$.


\end{thms}





\begin{notes}\label{r3.1.1} 


На основе сетки $\od_h$ построим триангуляцию области $\od$;


треугольные элементы, получаемые разбиением элементарных


четырехугольных элементов диагональю, имеют вершинами узлы из


$\od_h$ (напр., \cite{11}); элементы, примыкающие к


границе $\varGamma$, имеют криволинейные стороны. В случае разностной


схемы (\ref{eq.3.4}), (\ref{eq.3.9}) для функции $\ov{z}(x)$,


$x\in\od$, являющейся линейным интерполянтом $z(x)$ на треугольных


элементах, выполняется оценка, подобная оценке (\ref{eq.3.10}):


$$


%% \label{eq.3.12}


|u(x)-\ov{z}(x)|\leq M\sum_{s=1,2}(\eps+N_s^{-1})^{-2}N_s^{-2},


\quad x\in\od,


$$


неулучшаемая по вхождению величин $N_s$, $\eps$.


\end{notes}








\section{О построении $\eps$-равномерно


сходящихся схем\protect\\ на локально сгущающихся сетках}


\label{section.4}





Заметим, что сингулярная компонента решения краевой задачи


(\ref{eq.2.1}) экспоненциально убывает при удалении от границы


$\varGamma$. Сингулярная компонента при $r(x,\varGamma)\geq\si$ не превосходит


величины $M\de$, где $\de$ --- достаточно малое число, когда


$\si=m_1^{-1}\eps\ln\de^{-1}$, $m_1$ --- произвольное число из


интервала $(0,m_0)$, $m_0=m_{0(\ref{eq.7.6})}$. Невязка разностной


схемы на решении краевой задачи велика, однако лишь на этой


окрестности, являющейся достаточно узкой при малых значениях параметра


$\eps$ (см. \cite{1}, \cite{2}, \cite{4}). Это свойство разностной схемы будем


использовать при построении схем, сходящихся с возможно малым дефектом.


\medskip





{\ref{section.4}.1.} Имея в виду возможность использования схем


на достаточно произвольных локально сгущающихся сетках для решения


краевой задачи, будет удобно ввести сбалансированные сетки ---


сетки с произвольным распределением узлов (по $x_1$ и $x_2$),


однако, имеющие на $\od$ общее число узлов порядка $\Oh{N_1N_2}$,


т.\,е. такого же порядка, как и в случае равномерных базовых сеток.


Таким образом, объем вычислительной работы (пропорциональный числу


узлов, в которых требуется найти решение сеточной задачи) одного


порядка для сбалансированных и равномерных сеток.


Сбалансированные сетки, вообще говоря, не являются прямым


произведением сеток по $x_1$ и $x_2$. Подстраивание ориентации


шаблонов сеточных уравнений к границе $\varGamma$ не предполагается.


\medskip





{\ref{section.4}.2.} Рассмотрим краевую задачу (\ref{eq.2.1}) и


класс разностных схем на основе классических аппроксимаций задачи


в случае локально сгущающихся ``кусочно--равномерных'' сеток ---


сеток, равномерных как в ближайшей окрестности границы $\varGamma$, так и


вне несколько большей ее окрестности.


\medskip





{\ref{section.4}.2.1.} Для простоты предполагаем, что граница


$\varGamma$ проходит через начало координат и на множестве $D^0$ ---


$m^0$-окрестности начала координат, задается уравнением


$x_2=x_1$, причем, множеству $D$ принадлежат точки, для которых


выполняется соотношение $x_2<x_1$. Пусть каким--то образом на


$\od$ построена сетка


\begin{equation}


\label{eq.4.1}


\od{}_h^*=\od{}_h^*(D_1,D_2)=\od{}_h^*(\rho_1),


\end{equation}


где $\rho_1>0$ --- параметр, определяющий распределение узлов


сетки (\ref{eq.4.1}); этот параметр выбирается ниже. Эта сетка


является равномерной на множествах $D_1$ и $D_2$, $D_1,D_2\subset


\od{}^0$, где $D_1=D_1(\rho_1)\subset \od$ есть


$\rho_1$-окрестность границы $\varGamma$ (из $m^0$-окрестности начала


координат), $D_2=\{D\cap D^0\}\setminus\od_1(M\rho_1)$. Сетка


$D_{ih}=D_i\cap\od{}^*_{h(\ref{eq.4.1})}$ имеет шаг $h_{is}$ по


$x_s$, $s=1,2$. Считаем для простоты, что шаблоны схем, имеющие


центром узлы из $D_{ih}$, правильные, т.\,е. их ``плечи'' равны как


по направлению $x_1$, так и по $x_2$, и что выполнено условие


$$


a_{ss}(x)=c_0(x)\equiv 1,\quad a_{12}(x)=b_s(x)=c(x)\equiv 0,


\quad x\in\od, \quad s=1,2.


$$





Рассмотрим фрагменты сеточной задачи из класса разностных схем на


сетках (\ref{eq.4.1}), а именно, фрагменты на множествах $D_{1h}$


и $D_{2h}$. Пусть $z_i(x)$, $x\in\od_{ih}$, --- решение сеточной задачи


\begin{alignat}{2} \label{eq.4.2}


\La_{(\ref{eq.3.4})} z_i(x)&=f(x), &\quad x&\in D_{ih}; \\


z_i(x)&=u(x), &\quad x&\in \varGamma_{ih},\quad i=1,2, \notag


\end{alignat}


где $u(x)$, $x\in\od$, --- решение задачи (\ref{eq.2.1}).





Для функций $z_i(x)$, $x\in\od_{ih}$, с учетом априорных оценок


(\ref{eq.7.2}), (\ref{eq.7.6}) получаем оценки


\begin{subequations} \label{eq.4.3}


\begin{align} \label{eq.4.3a}


|u(x)-z_1(x)|&\leq M\sum_{s=1,2}(\eps+h_{1s})^{-2}h_{1s}^2,


\quad x\in \od_{1h}; \\


\label{eq.4.3b}


|u(x)-z_2(x)|&\leq M\Big\{ \Big[\sum_{s=1,2}


(\eps+h_{2s})^{-2}h_{2s}^2\Big]\max_{\varGamma_{2h}}


|V(x)|+\sum_{s=1,2}h_{2s}^2\Big\}, \ \ x\in \od_{2h},


\end{align}


\end{subequations}


где $V(x)$ --- сингулярная компонента решения краевой задачи


из представления (\ref{eq.7.1});


оценки (\ref{eq.4.3a}) и (\ref{eq.4.3b}) неулучшаемы по вхождению


величин $h_{1s}$, $\eps$ и $h_{2s}$, $\eps$ соответственно.


Эти оценки устанавливаются подобно оценке (\ref{eq.3.10}).





В силу оценки (\ref{eq.7.2b}) для того чтобы функция $z_2(x)$


сходилась $\eps$--равномерно, необходимо, чтобы для величины


$\rho_1$ выполнялось условие ($\rho_1\gg \eps$)


\begin{equation}\label{eq.4.4}


\eps=o(\rho_1), \quad \rho_1 \in (0,m].


\end{equation}


Оценка для функции $z_1(x)$, оптимальная по порядку сходимости


относительно $h_{1s}$, $s=1,2$, при фиксированном и


равном $MN_1N_2$ числе узлов сетки $\od_{1h}$


получается при условии $h_{11}\approx h_{12}$:


\begin{equation}\label{eq.4.5}


|u(x)-z_1(x)|\leq M\eps^{-2}\rho_1(N_1N_2)^{-1}


[1+\eps^{-2}\rho_1(N_1N_2)^{-1}]^{-1},\quad x\in \od_{1h};


\end{equation}


оценка неулучшаема по вхождению величин $\rho_1$, $(N_1N_2)^{-1}$ и


$\eps$. Из оценки (\ref{eq.4.5}) следует, что функция $z_1(x)$ не сходится


$\eps$-равномерно при $N_1,N_2\ra\infty$ при условии (\ref{eq.4.4}).





Таким образом, не существует кусочно-равномерных сеток


$\od{}_{h(\ref{eq.4.1})}^*$, на которых решения задач (\ref{eq.4.2})


при $i=1,2$ сходятся $\eps$-равномерно к решению задачи


(\ref{eq.2.1}). В случае вспомогательных задач (\ref{eq.4.2}) на


сетках (\ref{eq.4.1}) аналогично 


убеждаемся, что не существует сеток, на которых решения задач


(\ref{eq.4.2}) сходятся даже при условии


\begin{equation}\label{eq.4.6}


N_1^{-1}+N_2^{-1}\geq\eps^{1/2}.


\end{equation}





{\ref{section.4}.2.2.} Отсюда вытекает, что в случае семейства


сеток (\ref{eq.4.1}), как и семейства сеток


\begin{equation}\label{eq.4.7}


\od{}_h^*,


\end{equation}


являющихся равномерными в $\rho$-окрестности криволинейной границы $\varGamma$,


где $m\eps\leq\rho\leq M\eps$, и классических сеточных аппроксимаций


задачи, не существует схем, сходящихся при условии~(\ref{eq.4.6}).





Таким образом, справедлива





\begin{thms} \label{t4.1}


Для краевой задачи $(\ref{eq.2.1})$ в классе сбалансированных


разностных схем, строящихся на основе классических сеточных


аппроксимаций задачи на локально сгущающихся сетках $(\ref{eq.4.7})$,


не существует схем, дефект которых меньше, чем $2^{-1}$.


\end{thms}





\begin{notes} \label{r4.1.1}


Из приведенных построений вытекает, что в случае задачи


(\ref{eq.2.1}) использование локально сгущающихся сеток не


позволяет существенно ослабить условие (\ref{eq.3.11}) --- условие


сходимости классических разностных схем; на достаточно общих


локально сгущающихся сетках невозможно понизить порядок параметра


$\eps$ в условии (\ref{eq.2.3}) больше, чем в~2 раза. Это


утверждение согласуется с утверждением замечания \ref{r5.1.1} к


теореме \ref{t5.1} (имея в виду соотношение $P_{(\ref{eq.5.1})}=N_1N_2$).


\end{notes}








\section{О построении аппроксимаций решений задачи (\ref{eq.2.1})}


\label{section.5}





Рассмотрим некоторые проблемы, связанные с триангуляцией области


$\od$, которые возникают при аппроксимации решений сингулярно


возмущенной задачи (\ref{eq.2.1}). При этом используем


аналог поперечников по Колмогорову (см., напр., \cite{9}, \cite{10} и


библиографию там же).


\medskip





\ref{section.5}.1. Опишем аппроксимации множества


$\cu$ --- множества решений класса краевых задач (\ref{eq.2.1})


(определяемого условиями (\ref{eq.2.2})) в пространстве $X$ ---


множестве непрерывных функций с равномерной нормой. Решения задач


считаем достаточно гладкими на $\od$; для решений и их компонент


из представления (\ref{eq.7.1}) выполняются оценки


(\ref{eq.7.2}), (\ref{eq.7.6}).





Пусть $\od{}^h$ --- набор точек


(``сетка'') на $\od$. Сетки $\od{}^h$ могут быть как структурными


(порождаемыми некоторым регулярным семейством линий), так и


неструктурными. Число узлов сетки $\od{}^h$ обозначим через $P$.


Пусть $T_P$ есть триангуляция $\od$, порождаемая сеткой $\od{}^h$


(напр., \cite{11}). Для простоты считаем, что узлы сетки


$\od{}^h$ являются вершинами треугольных элементов, причем,


треугольные элементы образованы отрезками прямых, проходящих через


узлы $\od{}^h$. Пусть на множестве $\od{}^h$ определена некоторая


сеточная функция $u^h(x)$, $x\in\od{}^h$, через $\ov u{}^h(x)$,


$x\in\od$, обозначим ее линейный интерполянт. Совокупность таких


интерполянтов при фиксированной триангуляции $T_P$ обозначим через


$U^h_P$; совокупность всевозможных допустимых сеточных множеств


$\od{}^h$ и триангуляции $T_P$ на их основе обозначим через


$\cal{T}_P$. Эта совокупность $\cal{T}_P$ и совокупность


интерполянтов $U^h_P$ (для каждой триангуляции из $\cal{T}_P$)


определяют пространство $X$. Определим величину $d_P(\cu,X)$


(поперечник) соотношением


\begin{equation}\label{eq.5.1}


d_P(\cu,X)=\inf_{\cal{T}_P}\,\sup_{u\in\cu}\,


\inf_{\ov{u}^h\in U_P^h}\|u-\ov{u}^h\|,


\end{equation}


где $\|\cdot\|$ --- норма в $C$ (определение поперечника по Колмогорову


см., напр., в \cite{10}, гл.\,3).





Введем обозначения и определения.


Через $\rho_1(T_P^j)$ и $\rho_2(T_P^j)$, где


$T_P^j$ --- треугольный элемент из разбиения $T_P$, обозначим


радиусы вписанной и описанной окружностей для элемента $T_P^j$, $j


= 1,\dotsc,J$; здесь $J=J(P)$ --- число треугольных элементов


разбиения $T_P$ (считаем выполненным условие $J\approx P$).


Триангуляцию $T_P$ назовем {\it изотропной}, если выполняется условие


$$


\rho_1^{-1}(T_P^j)\rho_2(T_P^j)\leq M,\quad j=1,\dotsc,J,


$$


однако величины $\rho_1^{-1}(T_P^j)$, $\rho_2(T_P^j)$ могут


сильно отличаться от элемента к элементу, и {\it анизотропной} (с


коэффициентом анизотропии $\eta\geq M_0$, где $M_0$ может быть


достаточно большим), если выполняется условие


$$


\rho_1^{-1}(T_P^j)\rho_2(T_P^j)\leq M\eta,\quad j=1,\dotsc,J;


$$


постоянная $M$ не зависит от параметров $\eps$, $P$. Считаем, что


совокупность $\cal{T}_P$ определяется величиной $\eta$;


$d_P(\cu,X)=d_P(\cu,X;\eta)$. Подобным образом вводятся изотропная


и анизотропная триангуляции на подмножествах из $\od$. В том


случае, когда поперечники рассматриваются на множестве


$\od{}^0\subset\od$ (в этом случае поперечник обозначаем через


$d_P(\cu,X;\od{}^0)$), величина $\|u-\ov u{}^h\|$ в (\ref{eq.5.1})


вычисляется лишь по треугольным элементам, целиком принадлежащим


$\od{}^0$. Элементы триангуляции, для которых выполняется условие


$$


\eta\ra\infty, \ \text{ например, при } \ P\ra\infty


\ \text{ и/или } \ \eps\ra 0,


$$


назовем {\it существенно анизотропными}.





Величина $d_P(\cu,X)$ стремится к нулю при $P\ra\infty$, однако,


эта сходимость к нулю не~является $\eps$-равномерной.


Вообще говоря, приближения сходятся при $P\to\infty$ лишь при некоторых


соотношениях между величинами $P$ и $\eps$.





Пусть при $P\to\infty$ поперечник $d_P(\cu,X)$ (погрешность при оптимальной


точности приближения множества $\cu$ в пространстве $X$, или, короче,


погрешность оптимальной аппроксимации) стремится к нулю в случае


$P^{-1/2}=o(\eps^{\nu})$, $\eps\in E_{\eps}$, однако, сходимость при


условии $P^{-1/2}=\cal O(\eps^{\nu})$, вообще говоря, не имеет места.


В этом случае будем говорить, что поперечник {\it сходится} к нулю {\it с


дефектом} $\nu$.





Наибольший интерес представляют приближения, сходящиеся с возможно


меньшим дефектом и, в частности, сходящиеся $\eps$-равномерно.


\medskip





\ref{section.5}.2. Получим оценку поперечника,


когда граница $\varGamma$ на $\od{}^0_{(\ref{eq.4.1})}$ является


отрезком прямой.





Исследование модельных задач, подобное приведенному в


разделе~\ref{section.4}, показывает, что в случае изотропных


разбиений области $\od$ для поперечника выполняется оценка снизу


\begin{equation}


\label{eq.5.2}


d_P(\cu,X)\geq m(1+\eps P)^{-1}.


\end{equation}


На множестве


\begin{equation}\label{eq.5.3}


\od{}^1=\od_{1(\ref{eq.4.1})}(\rho_1),\quad \rho_1=M\eps,


\end{equation}


--- $\rho_1$-окрестности границы $\varGamma$ (из $\od{}^0_{(\ref{eq.4.1})}$),


получается оценка


\begin{equation}\label{eq.5.4}


d_P(\cu,X;\od{}^1_{(\ref{eq.5.3})})\geq m(1+\eps P)^{-1},


\end{equation}


неулучшаемая по вхождению величин $P$, $\eps$.





В случае анизотропных разбиений справедлива оценка


\begin{equation}\label{eq.5.5}


d_P(\cu,X)\geq m(1+\eps\eta P)^{-1}.


\end{equation}


На множестве $\od{}^1_{(\ref{eq.5.3})}$ имеем оценку


\begin{equation}\label{eq.5.6}


d_P(\cu,X;\od{}^1)\geq m(1+\eps\eta P)^{-1};


\end{equation}


оценка неулучшаема по вхождению величин $P$, $\eps$, $\eta$.





\begin{thms} \label{t5.1}


Пусть для компонент решений краевой задачи $(\ref{eq.2.1})$ из их


представлений $(\ref{eq.7.1})$, $(\ref{eq.7.5})$ выполняются


оценки $(\ref{eq.7.2})$, $(\ref{eq.7.6})$, где $K=2$. Тогда в


случае прямолинейной границы $\varGamma$ на $\od{}^0$ для поперечников


$d_P(\cu,X)$, $d_P(\cu,X;\od{}^1)$ на изотропной $($анизотропной$)$ 


триангуляции $T_P$ выполняются оценки


$(\ref{eq.5.2})$, $(\ref{eq.5.4})$ $($оценки $(\ref{eq.5.5})$,


$(\ref{eq.5.6}));$ оценки $(\ref{eq.5.4})$ и $(\ref{eq.5.6})$


неулучшаемы по вхождению величин $P$, $\eps$ и $P$, $\eps$,


$\eta$ соответственно.


\end{thms}





\begin{notes} \label{r5.1.1}


В силу оценки (\ref{eq.5.4}), а также оценки (\ref{eq.5.6}) при


условии, что величина $\eta$ не~зависит от $P$ и $\eps$,


неулучшаемый дефект сходимости погрешности оптимальной


аппроксимации на множестве $\od{}^1_{(\ref{eq.5.3})}$ не меньше, чем $2^{-1}$;


дефект сходимости поперечников $d_P(\cu,X)$ на изотропной и


анизотропной триангуляциях не меньше, чем $2^{-1}$.


\end{notes}





\begin{notes} \label{r5.1.2}


Условие


$$


\eta\ra\infty \ \ \text{при} \ \ P\ra\infty, \quad \eps\in E_{\eps},


\quad \eps P=\cal O(1)


$$


является необходимым и достаточным для того, чтобы дефект


сходимости на анизотропной триангуляции $T_P$ был меньше, чем $2^{-1}$.


Таким образом, использование существенно анизотропных триангуляций


является необходимым для того чтобы поперечник $d_P(\cu,X)$


сходился\linebreak $\eps$-равномерно либо с дефектом меньшим, чем $2^{-1}$.


\end{notes}





\begin{notes} \label{r5.1.3}


Условие


\begin{equation}\label{eq.5.7}


\eta=\eta(P,\eps)=\eps^{-1}\eta_0(P)


\end{equation}


является необходимым для $\eps$-равномерной ограниченности


поперечника, рассматриваемого на $\od{}^1$. При дополнительном


условии $\big(\eta_0(P)\,P \gg 1\big)$:


\begin{equation}\label{eq.5.8}


\eta_0^{-1}(P)=o(P), \quad P \geq M,


\end{equation}


имеем оценку сверху


$$


d_P(\cu,X;\od{}^1)\leq M\,P^{-1}\eta_0^{-1}(P),


$$


т.\,е. поперечник $d_P(\cu,\,X;\ \od{}^1)$ сходится $\eps$-равномерно.


Для поперечника $d_P(\cu,X;\od{}^0\cap\od)$ при условии (\ref{eq.5.7})


получается оценка


$$


d_P(\cu,X;\od{}^0\cap\od)\leq M[\eta_0^{-1}(P)P^{-1}\ln P+P^{-1/2}];


$$


при условии $\big(\eta_0^{-1}(P)\,P^{-1}\,\ln P \ll 1\big)$:


\begin{equation}\label{eq.5.9}


\eta_0^{-1}(P)=o(P\ln^{-1}P), \quad P \ra \infty,


\end{equation}


поперечник $d_P(\cu,X;\od{}^0\cap\od)$ сходится


$\eps$-равномерно. Таким образом, условия (\ref{eq.5.7}),


(\ref{eq.5.8}) являются необходимыми, а условия (\ref{eq.5.7}),


(\ref{eq.5.9}) --- достаточными для $\eps$-равномерной сходимости


(при $P\ra\infty$) поперечника $d_P(\cu,X;\od{}^0\cap\od)$ при


анизотропном разбиении.


\end{notes}





\ref{section.5}.3. Пусть множество $\varGamma$ --- криволинейная


граница из $\od{}^0_{(\ref{eq.4.1})}$ (с ограниченной кривизной) и


триангуляция $T_P$ является анизотропной (стороны


треугольных элементов --- отрезки прямых). Определим поперечник


$d_P^*(\cu,X)$ соотношением


\begin{equation}


\label{eq.5.10}


d_P^*(\cu,\,X)=\inf_{\eta}d_P(\cu,X;\eta).


\end{equation}





Справедлива





\begin{thms} \label{t5.2}


Пусть выполняется условие теоремы $\ref{t5.1}$. Тогда в случае


криволинейной границы $\varGamma$ на $\od{}^0_{(\ref{eq.4.1})}$ для


поперечника $d_{P(\ref{eq.5.10})}^*(\cu,X;\od{}^1)$, где


$\od{}^1=\od{}^1_{(\ref{eq.5.3})}$, выполняется оценка


$$


d_P^*(\cu,X; \od{}^1)\geq m(1+\eps^{1/2} P)^{-1},


$$


неулучшаемая по вхождению величин $P$, $\eps$.


\end{thms}





\begin{notes} \label{r5.2.1}


Дефект сходимости для поперечника $d_{P}^*(\cu,X;\od{}^1)$ (в


случае криволинейной границы $\varGamma$ на $\od{}^0$) равен $4^{-1}$. Таким


образом, на треугольных элементах с прямолинейными сторонами


дефект сходимости для поперечника $d_{P}^*(\cu,X)$ меньше, чем


$4^{-1}$, недостижим. Для построения аппроксимаций с дефектом


сходимости для поперечника $d_{P}^*(\cu,X)$ меньше, чем $4^{-1}$,


необходимо использовать треугольные элементы, являющиеся в


окрестности пограничного слоя криволинейными и существенно


анизотропными с локальной направленностью анизотропии


(направлением наибольшей стороны треугольных элементов),


согласованной с границей $\varGamma$.


\end{notes}





\begin{notes} \label{r5.2.2}


На множестве $\od$ в окрестности погранслоя можно ввести локальные


координаты, в которых множество $\varGamma$ становится прямолинейным, и


построить триангуляцию (на треугольных элементах с прямолинейными


сторонами), на которой поперечник сходится в окрестности


переходного слоя $\eps$-равномерно (см., напр., замечание


\ref{r5.1.3} к теореме \ref{t5.1}). В исходных переменных такая


триангуляция порождается криволинейными треугольными элементами.


На основе этой локальной триангуляции строится триангуляция


множества $\od$, на которой поперечник, рассматриваемый на $\od$,


сходится $\eps$-равномерно.


\end{notes}





Свойства поперечников, подобные приведенным в теореме \ref{t5.2} и


замечаниях \ref{r5.2.1}, \ref{r5.2.2}, сохраняются и в том случае,


когда для аппроксимации решений задачи (\ref{eq.2.1}) используются


интерполянты, приближающие решения задачи (\ref{eq.2.1}) с более


высоким порядком точности по сравнению с линейными интерполянтами.





Замечания \ref{r5.2.1}, \ref{r5.2.2} к теореме \ref{t5.2}


устанавливают необходимые условия (замечание \ref{r5.2.1}), при


которых погрешность оптимальной аппроксимации решения краевой


задачи слабо зависит от параметра $\eps$, и достаточные условия


(замечание \ref{r5.2.2}) для $\eps$-равномерной сходимости такой


погрешности. Требования, вытекающие из этих условий, полагаются в


основу при построении специальных схем.








\section{Схема на сетках, согласованных в погранслое с границей}


\label{section.6}





Чтобы продемонстрировать использование результатов


раздела~\ref{section.5}, приведем $\eps$-равномерно сходящуюся


разностную схему для задачи (\ref{eq.2.1}). При построении схемы


для краевой задачи (\ref{eq.2.1}) (с лучшим условием сходимости по


сравнению с (\ref{eq.3.11})) переформулируем задачу, перейдя (в


окрестности погранслоя) к связанным с границей $\varGamma$ переменным, в


которых граница является уже прямолинейной. Для задачи в новых


переменных построим разностную схему на прямоугольных сетках (в


частности, схему, сходящуюся $\eps$-равномерно) и затем вернемся


к старым переменным. В исходных переменных получающиеся сетки


(сетки, согласованные с границей) уже не будут прямоугольными,


что, вообще говоря, влечет некоторые неудобства при построении


сеточных областей и численном решении задачи. Однако в силу


результатов раздела~\ref{section.5}, использование


``криволинейных'' сеток, согласованных с границей $\varGamma$, является


необходимым. Рассмотрим разностную схему на основе такого подхода.


\medskip





\ref{section.6}.1. Построим разностную схему на основе


аппроксимации задачи (\ref{eq.7.4}), (\ref{eq.7.3}). Для простоты


считаем выполненным условие


\begin{equation}\label{eq.6.1^1}


r(x,\varGamma)=m_2,\quad x\in \varGamma^2,


\end{equation}


т.\,е. расстояние между границами $\varGamma$ и $\varGamma^2$ равно $m_2$;


$\od{}^2=D^2\cup\varGamma^2$.


На множестве $\ov{\wt{D}}{}^1$ введем прямоугольную сетку


\begin{subequations} \label{eq.6.1}


\begin{equation}\label{eq.6.1a}


\ov{\wt{D}}{}^1=\ov{\wt{\om}}_1\times\wt{\om}_2,


\end{equation}


где $\ov{\wt{\om}}_1$ и $\wt{\om}_2$ --- сетки на отрезке


$[0,m_1]$, $m_1=m_{1(\ref{eq.7.3})}$ (на оси $\xi_1$) и на оси


$\xi_2$ соответственно; сетка $\wt{\om}_2$ является периодической


(с периодом равным длине границы $\varGamma$). Пусть $\wt{N}_1+1$ и


$\wt{N}_2$ --- число узлов соответственно сеток $\ov{\wt{\om}}_1$


и $\wt{\om}_2$ (на интервале периодичности). На множестве


$\od{}^2$ строим сетку


\begin{equation}\label{eq.6.1b}


\od{}^2_h


\end{equation}


\end{subequations}


подобно построению сетки (\ref{eq.3.7}) на $\od$;


$\od{}^2_h=\od{}^2_h(R^2_{h(\ref{eq.3.1})})$, где


$R^2_{h(\ref{eq.3.1})}$ --- базовая сетка.





Приведем интерполянты сеточных функций, используемые при


построении схемы. Сеточные множества $\ov{\wt{D}}{}^1_h$ и


$\ov{D}{}^2_h$ определяют некоторые разбиения множеств


$\ov{\wt{D}}{}^1$ и $\ov{D}{}^2$ на треугольные элементы. Стороны


треугольных элементов, соединяющих узлы из $\ov{\wt{D}}{}^1_h$, а


также внутренние узлы из $D^2$, являются отрезками прямых (такие


элементы назовем прямолинейными); стороны элементов из $\od{}^2$,


примыкающих к границе $\varGamma^2$, могут быть криволинейными (такие


элементы назовем криволинейными). Пусть на множествах


$\ov{\wt{D}}{}^1_h$ и $\ov{D}{}^2_h$ определены некоторые сеточные


функции $Z(\xi)$ и $z(x)$ соответственно. На $\ov{\wt{D}}{}^1$


строим интерполянт $\ov{Z}(\xi)$ (линейный на треугольных


элементах). На множестве $\ov{D}{}^2$ строим интерполянт


$\wh{z}(x)$ следующим образом. На прямолинейных треугольных


элементах из $\ov{D}{}^2$ строим линейный интерполянт $\ov{z}(x)$ и


полагаем $\wh{z}(x)=\ov{z}(x)$. На криволинейных элементах


перейдем к переменным $\xi=(\xi_1,\xi_2)$ и по значениям функции


$z_{\xi}(\xi)=z(x(\xi))$ в вершинах элементов построим линейные


(по $\xi_1$, $\xi_2$) интерполянты $z^0(\xi)=\ov{(z_{\xi})}(\xi)$.


Положим $\wh{z}(x)=z^0_x(x)$. На отрезках прямых, являющихся


общими сторонами прямолинейных и криволинейных элементов, функция


$\wh{z}(x)$, вообще говоря, терпит разрыв. Для определенности


функцию $\wh{z}(x)$ на таких сторонах считаем равной полусумме


пределов из прилегающих треугольников. Интерполянт $\wh{z}(x)$,


$x\in\od{}^2$, построен.





Задачу (\ref{eq.7.4}), (\ref{eq.7.3}) аппроксимируем сеточной задачей


\begin{subequations} \label{eq.6.2}


\begin{align}\label{eq.6.2a}


\wt{\La} Z(\xi) & = F(\xi), \quad


\xi \in \wt{D}_h^1, \\


Z(\xi) & = \begin{cases}


\Phi(\xi), & \xi\in\wt{\varGamma}{}^1_{h}\cap \wt{\varGamma}; \\


(\wh{z})_{\xi}(\xi), & \xi\in\wt{\varGamma}{}^1_{h}\setminus \wt{\varGamma},


\end{cases}


\notag \\


\label{eq.6.2b}


\La z(x) & = f(x), \quad x\in D_h^2, \\


z(x) & = (\ov{Z})_x(x), \quad x\in\varGamma^2_{h}.\notag


\end{align}


\end{subequations}


Здесь $\La=\La_{(\ref{eq.3.3^1})}(L_{(\ref{eq.2.1})};\od{}^2_h)$;


оператор $\wt{\La}$ строится на сетке $\ov{\wt{D}}{}^1_h$ по


оператору $\wt{L}_{(\ref{eq.7.4})}$ подобно оператору


$\La_{(\ref{eq.3.3^1})}$


($\wt{\La}=\wt{\La}(\wt{L}_{(\ref{eq.7.4})};\ov{\wt{D}}{}^1_h)$,


причем $\eta(\xi)=1$ при $\xi\in \ov{\wt{D}}{}^1_h$);


$(\wh{z})_{\xi}(\xi)$ и $(\ov{Z})_x(x)$ суть соответственно


интерполянты $\wh{z}(x)$ и $\ov{Z}(\xi)$, переписанные в


переменных $\xi$ и $x$; $\wh{z}(x)$, $x\in\od{}^2$, и $\ov{Z}(\xi)$,


$\xi\in \ov{\wt{D}}{}^1$, --- интерполянты на треугольных элементах


разбиений областей $\od{}^2$ и $\ov{\wt{D}}{}^1$, порождаемых сетками


$\od{}^2_h$ и $\ov{\wt{D}}{}^1_h$ соответственно.





Функцию


$$


u^h(x)=\begin{cases}


(\ov{Z})_x(x), & x\in\{\ov{\wt{D}}{}^1\}_x; \\


\wh{z}(x),  & x\in\od\setminus \{\ov{\wt{D}}{}^1\}_x,


\end{cases}


\quad x \in \od,


$$


назовем решением разностной схемы (\ref{eq.6.2}), (\ref{eq.6.1})


--- схемы метода декомпозиции области в случае перекрывающихся


подобластей.





Операторы $\wt{\La}$ и $\La$ являются монотонными (на множествах


$\wt D{}^1_h$ и $D^2_h$) в том случае, когда выполняются


соответственно условия


\begin{subequations} \label{eq.6.3}


\begin{gather}\label{eq.6.3a}


\min_{s,\wt{D}_h^1} \{2(\wt{h}_s^{i_s}+


\wt{h}_s^{i_s-1})^{-1}A_{ss}(\xi)-


\max [(\wt{h}_{3-s}^{i_{3-s}})^{-1},


(\wt{h}_{3-s}^{i_{3-s}-1})^{-1}]


|A_{12}(\xi)|-|B_s(\xi)| \}\geq 0; \\


%


\label{eq.6.3b}


\min_{s,D_h^{2(in)}} \{2(h_s^{i_s}+h_s^{i_s-1})^{-1}a_{ss}(x)-


\max [(h_{3-s}^{i_{3-s}})^{-1},


(h_{3-s}^{i_{3-s}-1})^{-1} ]


|a_{12}(x)|-|b_s(x)| \}\geq 0.


\end{gather}


\end{subequations}


Здесь $D_h^{2(in)}$ --- множество строго внутренних узлов


(определяемых для множества $D_h^2$), строящихся подобно множеству


$D_{h(\ref{eq.3.3})}^{(in)}$; $\wt h{}_1^i$ и $\wt h{}_2^j$ --- шаги


сеток $\ov{\wt{\om}}_{1(\ref{eq.6.1})}$ и


$\wt{\om}_{2(\ref{eq.6.1})}$ соответственно.





Заметим, что условия (\ref{eq.6.3}) выполняются, например, когда


коэффициенты при смешанных производных равны нулю


$$


A_{12}(\xi)\equiv 0,\ \ \xi\in \ov{\wt{D}}{}^1,


\quad a_{12}(x)\equiv 0,\ \ x\in\od{}^2,


$$


либо когда эти коэффициенты достаточно малы (по сравнению с


коэффициентами $A_{ss}(\xi)$, $a_{ss}(x)$) и шаги сеток $\om_1$,


$\om_2$, а также $\ov{\wt{\om}}_1$, $\wt{\om}_2$ одного порядка;


такое условие имеет место, например, в случае задачи (\ref{eq.2.1})


на круге при условии $L_2= \triangle$ в окрестности


границы~$\varGamma$ и малости коэффициента $a_{12}(x)$ по


сравнению с $a_{11}(x)$, $a_{22}(x)$ вне этой окрестности.





В случае условий (\ref{eq.6.3}) разностная схема (\ref{eq.6.2}),


(\ref{eq.6.1}) является монотонной. В том случае, когда условия


(\ref{eq.6.3}) нарушаются, при построении монотонных сеточных


аппроксимаций задачи (\ref{eq.2.1}) можно воспользоваться техникой


работы \cite{1}.





Для решения задачи (\ref{eq.2.1}) используем разностную схему


(\ref{eq.6.2}), (\ref{eq.6.1}) на сетках, сгущающихся в


пограничном слое. Сеточную задачу рассматриваем на сетках


\begin{subequations} \label{eq.6.4}


\begin{equation}\label{eq.6.4a}


\ov{\wt{D}}{}^1_h,\quad \od{}^2_h=\od{}^2_h(R^2_h),


\end{equation}


где $\ov{\wt{D}}{}^1_h$, $R^2_h$ --- равномерные сетки. Для простоты


считаем выполненным условие


\begin{equation}\label{eq.6.4b}


N_s=\wt{N}_s=N,\quad s=1,2.


\end{equation}


Параметры подобластей $\od{}^k$ определим соотношениями


\begin{gather} \label{eq.6.4c}


m_{2(\ref{eq.6.1^1})}=\si,\quad m_{1(\ref{eq.7.3})}=\si+\de, \\


%


\si=\si(\eps,N)=\min[2^{-1}r^0,2m_3^{-1}\eps\ln N],


\quad \de=\de(\eps)=m_4\eps,\notag


\end{gather}


\end{subequations}


где $r^0$ --- минимальный радиус кривизны границы $\varGamma$,


$m_3=m_{(\ref{eq.7.6})}$.





С учетом априорных оценок решений задачи (\ref{eq.2.1}) находим


оценку для решения схемы (\ref{eq.6.2}), (\ref{eq.6.4})


\begin{equation}\label{eq.6.5}


|u(x)-u^h(x)|\leq M N^{-2}\ln^2N,\quad x\in\od.


\end{equation}


Tаким образом, разностная схема сходится $\eps$-равномерно со


скоростью сходимости $\cal O(N^{-2}\ln^2N)$.





\begin{thms} \label{t6.1}


Пусть для решения краевой задачи $(\ref{eq.2.1})$ выполняются


условие теоремы $\ref{t3.1}$ и условия


$(\ref{eq.6.3})$. Тогда решение разностной схемы $(\ref{eq.6.2})$,


$(\ref{eq.6.4})$ при $N\ra \infty$ сходится к решению краевой


задачи $\eps$-равномерно. Для решения разностной схемы


справедлива оценка $(\ref{eq.6.5})$.


\end{thms}








Автор выражает глубокую признательность Н.С.\,Бахвалову за


внимание к исследованиям, приведенным в работе,


и предложение использовать подход на


основе поперечников по Колмогорову.








\begin{thebibliography}{99}








\bibitem{1}


Шишкин Г.И. {\em Сеточные аппроксимации сингулярно


возмущенных эллиптических и параболических уравнений}. --


Екатеринбург: УрО РАН, 1992. -- 233~c.


\bibitem{2}


Miller J.J.H., O'Riordan E., Shishkin G.I. {\em Fitted


numerical methods for singular perturbation problems$:$ error


estimates in the maximum norm for linear problems in one and two


dimensions}. -- Singapore: World Scientific, 1996. -- 166~p.


\bibitem{3}


Roos H.-G., Stynes M., Tobiska L. {\em Numerical methods for


singularly perturbed differential equations. Convection-diffusion


and flow problems}. -- Berlin: Springer, 1996. -- 348~p.


\bibitem{4}


Farrell P.A., Hegarty A.F., Miller J.J.H., O'Riordan E.,


Shishkin G.I. {\em Robust computational techniques for boundary


layers}. --- Boca Raton: Chapman \& Hall/CRC, 2000. -- 254~p.


\bibitem{5}


Shishkin G.I., Vabishchevich P.N. {\em Interpolation finite difference


schemes on grids locally refined in time} // Computer Methods in Applied


Mechanics and Engineering. -- 2000. -- V.\,190. -- \No\,8--10.


-- P.\,889--901.


\bibitem{6}


Марчук Г.И., Шайдуров В.В. {\em Повышение точности решений разностных схем}.


-- М.: Наука, 1979. -- 318~с.


\bibitem{7}


Birkhoff G., Lynch R.E. {\em Numerical solution of elliptic problems}. --


Philadelphia: SIAM, 1984. -- 319~p.


\bibitem{12}


Шишкин Г.И. {\em Численные методы на адаптивных сетках для сингулярно


возмущенных эллиптических уравнений в области с криволинейной границей} //


Изв. вузов. Математика. -- 2003. -- \No\,1. -- C.\,74--85.


\bibitem{13}


Ладыженская О.А., Уральцева Н.Н. {\em Линейные и квазилинейные уравнения 


эллиптического типа}. -- М.: Наука, 1973. -- 576~с.


\bibitem{8}


Самарский А.А. {\em Теория разностных схем}. -- М.: Наука, 1989. -- 616~c.


\bibitem{11}


Марчук Г.И. {\em Методы вычислительной математики}. -- М.: Наука, 1989. --


608~с.


\bibitem{9}


Бахвалов Н.С. {\em Численные методы}. -- М.: Наука, 1973. -- 632~с.


\bibitem{10}


Бабенко К.И. {\em Основы численного анализа}. -- М.: Наука, 1986. -- 744~с.





\end{thebibliography}





\vskip 2cm





\post{\it Институт математики и механики}{\it Поступила}


\post{\it Уральского отделения}{22.04.2004}


\post{\it Российской Академии наук}





\label{end}


\end{document}








620219 Екатеринбург, ГСП-384,\\


ул. С. Ковалевской, 16,\\


Институт математики и механики УрО РАН


shishkin@imm.uran.ru








{\sl О содержании работы.~} Разностные схемы на основе классических


аппроксимаций задачи (\ref{eq.2.1}) при использовании прямоугольных


сеток, являющихся прямым произведением сеток по $x_1$ и $x_2$,


рассматриваются в разделе 4. В разделе 5 обсуждаются проблемы,


связанные с построением схем, сходящихся при условии, более слабом,


чем условие (\ref{eq.2.3}), при использовании локально сгущающихся


прямоугольных сеток. Построение аппроксимаций решений задачи (\ref{eq.2.1})


на основе линейных интерполянтов, строящихся на триангуляциях области,


изучается в разделе 6. В разделе 7 на сетках, согласованных в погранслое с


границей области, строится схема, сходящаяся $\eps$-равномерно.


Априорные оценки, используемые в построениях, приводятся в разделе 3.


Читатель, заинтересованный лишь в оценках, которые понадобятся в дальнейшем для


обоснования строящихся схем, может пропустить вывод оценок и перейти сразу


к утверждению теоремы 3.1.


